%--------------------------------------------------------------------------------------
%--------------------------------------------------------------------------------------
%--------------------------------------------------------------------------------------
\section{Disponibilidad de datos}
\label{sec:datos-apendice}

Los insumos utilizados en la estimación son tres archivos: la fuente consolidada de covariables municipales de 2018 y 2024 (\texttt{\detokenize{Covariables_18&24.csv}}), con las variables y fuentes descritas en la Tabla~\ref{tab:matriz-auxiliar-ipm}; la estimación directa departamental del IPM 2024 con su varianza de muestreo (\texttt{\detokenize{IPM_Depto_2024.xlsx}}), descrita en la Sección~\ref{sec:datos}; y la cartografía municipal del Marco Geoestadístico Nacional del DANE (\texttt{\detokenize{MGN_ADM_MPIO_GRAFICO.shp}}), utilizada para construir la matriz de vecindad del Modelo 3. Las fuentes primarias de las covariables son de acceso público.
% La disponibilidad de archivos de entrada y restricciones de redistribucion se declarara en el apartado de repositorio de A.3.

%--------------------------------------------------------------------------------------
%--------------------------------------------------------------------------------------
%--------------------------------------------------------------------------------------
\section{Notación de variables codificadas}
\label{sec:tabla-variables-codificadas}


% Diccionario de variables - Apéndice A.2
\begingroup
\small
\renewcommand{\arraystretch}{1.1}
\setlength{\tabcolsep}{5pt}

\begin{longtable}{p{5cm} p{7cm}}

\caption{Correspondencia entre los nombres de variables utilizados en el procesamiento de datos y su descripción}
\label{tab:diccionario-variables}\\

\hline
\textbf{Nombre en código} & \textbf{Nombre descriptivo} \\
\hline
\endfirsthead

\multicolumn{2}{c}{\small\textit{Tabla \thetable\ -- continuación}}\\
\hline
\textbf{Nombre en código} & \textbf{Nombre descriptivo} \\
\hline
\endhead

\hline
\multicolumn{2}{r}{\small\textit{Continúa en la siguiente página}}\\
\endfoot

\hline
\endlastfoot

\texttt{\detokenize{cod_depto}} & Código del departamento \\
\texttt{\detokenize{cod_mpio}} & Código del municipio \\
\texttt{depto} & Nombre del departamento \\
\texttt{mpio} & Nombre del municipio \\

\texttt{\detokenize{pct_subsidiado_2024}} & Porcentaje de afiliados al régimen subsidiado (2024) \\
\texttt{\detokenize{pct_cont_2024}} & Porcentaje de afiliados al régimen contributivo (2024) \\
\texttt{\detokenize{aseguramiento_capado}} & Aseguramiento en salud truncado en 1 (2024) \\
\texttt{\detokenize{ratio_subs_cont_2024}} & Razón subsidiado/contributivo (2024) \\
\texttt{\detokenize{log_total_afiliados_2024}} & Logaritmo del total de afiliados en salud (2024) \\
\texttt{\detokenize{mortalidad_infantil}} & Tasa de mortalidad infantil \\
\texttt{\detokenize{log_mortalidad_infantil}} & Transformación logarítmica de la mortalidad infantil \\

\texttt{\detokenize{tasa_matriculacion_5_16}} & Tasa de matriculación de 5 a 16 años \\
\texttt{\detokenize{cobertura_neta}} & Cobertura neta educativa total \\
\texttt{\detokenize{cobertura_neta_media}} & Cobertura neta en educación media \\
\texttt{desercion} & Tasa de deserción escolar \\
\texttt{repitencia} & Tasa de repitencia escolar \\
\texttt{reprobacion} & Tasa de reprobación escolar \\

\texttt{\detokenize{Cobertura_acueducto}} & Cobertura de acueducto \\
\texttt{\detokenize{Cobertura_alcantarillado}} & Cobertura de alcantarillado \\
\texttt{\detokenize{Cobertura_aseo}} & Cobertura de aseo \\
\texttt{\detokenize{infraestructura_basica}} & Índice de infraestructura básica \\
\texttt{\detokenize{ICEE_rural}}, \texttt{\detokenize{cobertura_energia_rural}} & Cobertura de energía eléctrica rural \\

\texttt{\detokenize{viirs_rad}} & Luminosidad nocturna (radiancia VIIRS) \\
\texttt{\detokenize{log_viirs}} & Logaritmo de uno más la radiancia nocturna \\

\texttt{poblacion} & Población municipal 2024 \\
\texttt{\detokenize{poblacion_5_a_16}} & Población municipal de 5 a 16 años \\
\texttt{\detokenize{nacidos_vivos}} & Nacidos vivos \\
\texttt{\detokenize{pob_asegurada}} & Población afiliada al sistema de salud \\

\texttt{\detokenize{Recaudo_ICA}} & Recaudo total del impuesto de industria y comercio \\
\texttt{\detokenize{Rec_ICA_PC}} & Recaudo del impuesto de industria y comercio per cápita \\
\texttt{\detokenize{log_rec_ica_pc}} & Transformación logarítmica del recaudo per cápita \\

\texttt{\detokenize{ipm_directo}} & IPM directo departamental \\
\texttt{vardir} & Varianza directa de muestreo \\

\end{longtable}
\endgroup


%--------------------------------------------------------------------------------------
%--------------------------------------------------------------------------------------
%--------------------------------------------------------------------------------------
% Asegurar que la tabla de A.2 se coloque antes de iniciar A.3.
\clearpage

\section{Disponibilidad de códigos y reproducibilidad}
\label{sec:parte1}\label{sec:datos-reprocucibilidad}

Los resultados de esta investigación se obtuvieron mediante un programa principal desarrollado en R, denominado \texttt{modelamiento\_SAE\_IPM\_2024\_\allowbreak FINAL\_UNICO.R}, en el cual se integran de manera secuencial las etapas de preparación de datos, estimación de los modelos, validación, análisis de sensibilidad, robustez y exportación de resultados. Esta organización permite ejecutar los procedimientos metodológicos desde un único archivo, conservando los objetos y parámetros necesarios para cada etapa.

El programa comprende los siguientes componentes:

\begin{itemize}

    \item \textbf{Preparación y control de calidad de datos:} lectura e integración de las fuentes auxiliares municipales y de las estimaciones directas departamentales del IPM 2024; verificación de identificadores territoriales, normalización de unidades, controles de integridad, tratamiento de registros incompletos y construcción de las bases analíticas.

    \item \textbf{Estimación de los modelos M1, M2 y M3:} ajuste del modelo Fay--Herriot departamental mediante REML (M1), estimación de la formulación de desagregación municipal no lineal acotada (M2) y estimación del modelo con componente espacial de rango reducido y penalización inspirada en CAR (M3).

    \item \textbf{Validación y comparación de modelos:} ejecución de las pruebas de validez y estabilidad numérica, validación por exclusión de dominio (LODO) para los 33 departamentos, cálculo de RMSE, MAE y correlaciones, y comparación de las formulaciones bajo los escenarios de covariables E1, E2 y E3.

    \item \textbf{Benchmarking departamental:} aplicación de desplazamientos en escala logit para garantizar la coherencia de las estimaciones municipales agregadas con las estimaciones directas departamentales, así como evaluación de su magnitud y efectos.

    \item \textbf{Análisis y diagnósticos espaciales:} construcción de la matriz de vecindad municipal por contigüidad Reina, obtención de la base espacial de rango reducido, evaluación de la autocorrelación espacial y análisis de sensibilidad del Modelo 3 frente a diferentes valores de los parámetros espaciales.

    \item \textbf{Análisis de sensibilidad y robustez:} evaluación LODO de M3 para los valores de $\rho\in\{0.25,0.50,0.75,0.90\}$; identificación de departamentos potencialmente influyentes mediante diagnósticos OLS; exclusión individual de San Andrés y Providencia, Vichada y Vaupés; reestimación de M1, M2 y M3 bajo el escenario E2; y comparación de los coeficientes y métricas predictivas.

    \item \textbf{Generación y exportación de resultados:} consolidación de estimaciones municipales, indicadores de ajuste, métricas de validación, tablas comparativas, diagnósticos de robustez, archivos CSV, libros Excel y productos cartográficos utilizados en el documento.

\end{itemize}

Los análisis complementarios se ejecutan dentro del mismo programa y utilizan los objetos generados previamente, sin sustituir las estimaciones definitivas de los tres modelos. La reproducción completa requiere disponer de las fuentes de información descritas en la Sección~\ref{sec:datos-apendice} y configurar las rutas de entrada y salida del programa de acuerdo con el entorno de ejecución.


La ejecución documentada utilizó R 4.3.3 con los paquetes sae 1.3, spdep 1.3.1, spatialreg 1.3.1, sf 1.0.15, Matrix 1.6.5, RSpectra 0.16.1, dplyr 1.1.4, readxl 1.4.3, openxlsx 4.2.5.2, car 3.1.2, lmtest 0.9.40, numDeriv 2016.8.1.1 y ggplot2 3.4.4. Los parámetros de ejecución fueron los siguientes:

\begin{itemize}
    \item Semilla global 2026. M2: 12 arranques en el ajuste completo y 5 por pliegue LODO, con semilla 5000 más el índice del departamento excluido. M3: 8 arranques en el ajuste completo y 3 por pliegue, con semilla 12000 más 100 veces el índice del escenario más el índice del departamento excluido. Los análisis de sensibilidad utilizan semillas fijas documentadas en el programa unificado.
    \item Optimización de M2 y M3 con Nelder--Mead seguido de BFGS; tolerancia relativa de $10^{-10}$ y 15{,}000 iteraciones máximas en la primera etapa de M2, y de $10^{-9}$ con 12{,}000 y 5{,}000 iteraciones en M3. Se considera que dos arranques alcanzan el mismo mínimo cuando sus funciones objetivo difieren en menos de $10^{-4}$.
    \item Benchmarking por búsqueda de raíz en $[-25,25]$ de la escala logit, con tolerancia de $10^{-12}$. M1 se ajusta por REML mediante la función \texttt{eblupFH} del paquete sae.
    \item Vecindad por contigüidad Reina sobre cartografía del Marco Geoestadístico Nacional, con tolerancia de unión de 0.001, matriz binaria y normalización simétrica $\mathbf{S}=\mathbf{D}^{-1/2}\mathbf{W}\mathbf{D}^{-1/2}$; autovectores obtenidos con RSpectra. Se emplean $K=6$ y $\rho=0.50$ como referencia, con grillas $K\in\{3,4,6,8,10\}$ y $\rho\in\{0.25,0.50,0.75,0.90\}$.
    \item Índice de Moran departamental con 9{,}999 permutaciones y municipal con 999.
    \item En el análisis adicional de robustez E2, M2 utiliza 12 arranques en cada ajuste de 32 departamentos y 5 por pliegue de 31; M3 utiliza 8 y 3 arranques, respectivamente, con $K=6$ y $\rho=0.50$ fijos. Las semillas se fijan explícitamente en el programa unificado.
\end{itemize}

\paragraph{Salidas del análisis de robustez departamental.}

El bloque de análisis de robustez departamental, integrado en el programa principal \texttt{modelamiento\_SAE\_IPM\_2024\_\allowbreak FINAL\_UNICO.R}, genera los siguientes archivos en la carpeta relativa \texttt{salidas\_departamentos\_influyentes/\allowbreak robustez\_E2/}:

\begin{itemize}
    \item \texttt{01\_diagnosticos\_formales\_\allowbreak OLS\_E2.csv}: diagnósticos OLS de los 33 departamentos (Cook, apalancamiento, residuo estandarizado, máximo DFBETAS absoluto e indicadores de señal). Sustenta la Sección~\ref{sec:res-ols}.
    \item \texttt{02\_robustez\_metricas\_\allowbreak LODO\_E2.csv}: métricas RMSE, MAE y Pearson originales y reestimadas para cada exclusión y modelo, junto con los cambios de RMSE y MAE.
    \item \texttt{03\_robustez\_\allowbreak coeficientes\_E2.csv}: coeficientes originales y reestimados de M1, M2 y M3, diferencia y variación relativa porcentual. Las covariables de M2 y M3 se comparan en una escala común; los coeficientes espaciales se identifican separadamente.
    \item \texttt{04\_robustez\_\allowbreak predicciones\_E2.csv}: IPM directo, predicciones LODO originales y predicciones LODO reestimadas por modelo y departamento de evaluación.
\end{itemize}

Los tres departamentos seleccionados fueron San Andrés y Providencia, Vichada y Vaupés. Los cuatro archivos permiten reconstruir las comparaciones de influencia y sensibilidad presentadas en la Sección~\ref{sec:res-comparacion}, sin confundir la exclusión de una observación del cálculo del error con la reestimación efectiva de los modelos. Las métricas basales y reestimadas se calculan sobre los mismos 32 departamentos de evaluación, aunque las primeras proceden de pliegues entrenados con 32 dominios y las segundas con 31.

En un equipo de dos núcleos, la etapa de modelamiento principal registró un tiempo aproximado de 22 minutos y el análisis complementario de sensibilidad a $K$ cerca de 32 minutos. Estos tiempos corresponden a ejecuciones documentadas de las etapas respectivas y no representan el tiempo total de ejecución del programa unificado, que incluye análisis adicionales de sensibilidad y robustez.

\paragraph{Repositorio de reproducibilidad y acceso a materiales.}

El código y las salidas analíticas se organizarán en una carpeta digital de respaldo para facilitar la trazabilidad y la reproducción de los resultados de este trabajo. La ubicación definitiva del repositorio se incorporará en el siguiente campo antes de la entrega final:

\begin{quote}
\noindent\textbf{Enlace del repositorio:} \texttt{\detokenize{[https://github.com/BrySB/TG_MCD_Grupo5_SAE_Pobreza_2026]}}
\end{quote}

% AL PUBLICAR EL REPOSITORIO: sustituir la línea anterior por
% \noindent\textbf{Enlace del repositorio:} \url{https://enlace-real-del-repositorio}
% Confirmar antes de entregar: acceso de lectura, archivos efectivamente incluidos y permisos de redistribucion.

El material previsto comprende el programa unificado
\texttt{modelamiento\_SAE\_IPM\_2024\_\allowbreak FINAL\_UNICO.R}, los archivos de resultados consolidados y los cuatro archivos CSV de robustez identificados anteriormente. Los insumos de entrada y la cartografía se describen en la Sección~\ref{sec:datos-apendice}; su disponibilidad para descarga dependerá de las condiciones de acceso y redistribución correspondientes. Cuando un archivo no pueda compartirse, el repositorio deberá indicar la fuente oficial y las instrucciones para obtenerlo.

Para reproducir el procedimiento se requiere configurar las rutas de los datos, disponer de las dependencias y versiones de software señaladas en esta sección y ejecutar el programa unificado en el orden en que están definidos sus bloques. La publicación del repositorio y la verificación de su contenido quedan pendientes hasta asignar el enlace definitivo.

%--------------------------------------------------------------------------------------
\section{Errores LODO por departamento}
\label{sec:apendice-errores}

La Tabla~\ref{tab:errores-lodo-dep} presenta el error LODO de cada departamento en las nueve combinaciones de modelo y escenario, calculado como la estimación LODO menos la estimación directa; la columna IPM corresponde a la estimación directa. Valores positivos indican sobreestimación.

\setlength{\tabcolsep}{2pt}
\footnotesize
\begin{longtable}{l c c c c c c c c c c}
\caption{Error LODO por departamento, en puntos del IPM.}
\label{tab:errores-lodo-dep}\\
\hline
 & & \multicolumn{3}{c}{E1} & \multicolumn{3}{c}{E2} & \multicolumn{3}{c}{E3} \\
Departamento & IPM & M1 & M2 & M3 & M1 & M2 & M3 & M1 & M2 & M3 \\
\hline
\endfirsthead
\hline
Departamento & IPM & M1 & M2 & M3 & M1 & M2 & M3 & M1 & M2 & M3 \\
\hline
\endhead
\hline
\endfoot
Antioquia & 10.90 & 0.72 & $-1.22$ & 4.89 & 0.92 & $-1.11$ & 3.49 & 1.91 & 0.05 & 3.70 \\
Atlántico & 9.50 & $-9.20$ & $-4.84$ & $-5.86$ & $-9.01$ & $-4.68$ & $-5.38$ & $-7.45$ & $-3.43$ & $-5.39$ \\
Bogotá D.C. & 5.40 & $-6.56$ & $-2.61$ & $-3.62$ & $-5.55$ & $-1.66$ & $-3.59$ & $-0.39$ & 1.25 & $-2.66$ \\
Bolívar & 15.60 & 1.79 & $-0.06$ & 1.37 & $-1.31$ & $-1.62$ & 2.29 & $-1.67$ & $-1.69$ & 0.63 \\
Boyacá & 6.90 & 1.34 & 2.85 & 1.92 & 2.23 & 3.37 & 1.71 & 1.09 & 2.71 & 4.36 \\
Caldas & 9.20 & $-1.27$ & $-1.26$ & $-0.19$ & 2.71 & 1.35 & 3.23 & 3.84 & 2.85 & 4.73 \\
Caquetá & 13.90 & 4.58 & 5.16 & 4.02 & 6.84 & 6.47 & 1.44 & 6.45 & 6.19 & 0.67 \\
Cauca & 14.00 & 0.62 & 0.22 & $-0.94$ & 3.35 & 2.74 & $-0.96$ & 1.11 & $-0.02$ & $-5.49$ \\
Cesar & 13.40 & $-0.19$ & $-3.06$ & $-1.63$ & 0.32 & $-1.17$ & 2.62 & 7.41 & 5.80 & 8.17 \\
Córdoba & 25.70 & $-7.34$ & $-10.70$ & $-6.99$ & $-8.04$ & $-10.36$ & $-4.45$ & $-6.61$ & $-8.61$ & $-0.30$ \\
Cundinamarca & 7.40 & 11.90 & 5.74 & 2.29 & 10.21 & 3.95 & 0.66 & 9.48 & 3.85 & 1.55 \\
Chocó & 33.90 & $-1.88$ & 1.68 & 13.37 & $-10.61$ & $-12.54$ & $-7.46$ & $-12.10$ & $-15.86$ & $-11.98$ \\
Huila & 10.90 & $-3.93$ & $-0.65$ & $-2.96$ & $-2.79$ & $-0.77$ & $-4.42$ & $-3.34$ & $-0.85$ & $-4.54$ \\
La Guajira & 39.30 & $-2.89$ & $-7.37$ & $-10.48$ & $-12.91$ & $-15.65$ & $-13.33$ & $-11.54$ & $-14.44$ & $-11.46$ \\
Magdalena & 19.20 & 5.04 & 4.91 & 12.36 & 2.17 & 1.72 & 12.15 & 0.24 & $-0.51$ & 4.83 \\
Meta & 10.40 & 6.11 & 5.24 & 1.87 & 4.72 & 3.49 & $-0.43$ & 3.49 & 2.80 & 1.95 \\
Nariño & 18.10 & $-3.86$ & $-3.08$ & $-4.76$ & 0.80 & 0.90 & $-6.55$ & $-0.80$ & $-1.20$ & $-12.39$ \\
N. de Santander & 15.20 & $-6.15$ & $-5.01$ & $-2.11$ & $-5.14$ & $-4.39$ & $-2.42$ & $-5.72$ & $-4.93$ & $-5.32$ \\
Quindío & 7.40 & $-2.21$ & $-0.76$ & $-2.10$ & $-2.69$ & $-1.13$ & $-1.96$ & 1.10 & 1.79 & $-0.89$ \\
Risaralda & 9.50 & 1.30 & $-0.32$ & $-0.61$ & $-1.19$ & $-1.14$ & $-0.84$ & $-1.58$ & $-0.70$ & $-1.86$ \\
Santander & 6.80 & 1.55 & 2.39 & 4.56 & 1.30 & 1.84 & 2.89 & 0.97 & 1.74 & 2.98 \\
Sucre & 21.80 & $-2.58$ & $-5.10$ & 0.66 & $-4.19$ & $-5.58$ & 1.26 & $-3.44$ & $-4.90$ & 1.17 \\
Tolima & 12.60 & $-0.71$ & $-0.43$ & $-3.35$ & $-3.55$ & $-2.28$ & $-3.79$ & $-3.48$ & $-2.56$ & $-3.68$ \\
Valle del Cauca & 6.20 & $-0.10$ & 0.71 & 0.59 & 4.71 & 3.87 & 2.74 & 6.16 & 5.56 & 2.59 \\
Arauca & 17.70 & $-2.27$ & $-1.84$ & $-2.94$ & $-2.94$ & $-2.42$ & $-4.09$ & $-4.61$ & $-4.14$ & $-2.31$ \\
Casanare & 11.30 & $-1.58$ & $-0.35$ & $-3.26$ & $-2.51$ & $-0.73$ & $-3.77$ & $-2.28$ & $-0.71$ & $-0.56$ \\
Putumayo & 11.80 & 4.31 & 2.79 & $-1.03$ & 6.87 & 7.51 & 0.71 & 7.72 & 6.81 & $-4.57$ \\
San Andrés & 6.00 & 18.11 & 10.82 & 7.36 & 19.34 & 13.54 & 6.04 & 2.00 & 1.87 & $-2.01$ \\
Amazonas & 15.60 & 17.67 & 15.96 & 12.12 & 15.82 & 13.20 & 7.12 & 19.86 & 22.17 & 18.65 \\
Guainía & 49.00 & $-10.79$ & $-4.66$ & $-7.11$ & $-8.69$ & 2.18 & $-1.12$ & $-15.18$ & $-15.76$ & $-17.31$ \\
Guaviare & 21.30 & $-2.56$ & 1.54 & $-0.78$ & 4.90 & 6.64 & 1.37 & 1.68 & 1.37 & $-0.94$ \\
Vaupés & 37.40 & 10.23 & $-2.37$ & 3.28 & 13.85 & 9.85 & 4.15 & 17.89 & 19.33 & 33.31 \\
Vichada & 70.20 & $-23.54$ & $-6.53$ & $-15.86$ & $-15.80$ & $-9.87$ & $-2.79$ & $-17.37$ & $-19.83$ & $-32.70$ \\
\end{longtable}
\normalsize
\setlength{\tabcolsep}{6pt}

%--------------------------------------------------------------------------------------
\section{Desplazamientos de benchmarking por departamento}
\label{sec:apendice-benchmarking}

La Tabla~\ref{tab:delta-dep} presenta el desplazamiento $\hat{\delta}_d$ de la Ecuación~\ref{eq:benchmarking}, en la escala logit, para cada departamento, modelo y escenario.

\footnotesize
\begin{longtable}{l c c c c c c}
\caption{Desplazamiento de benchmarking $\hat{\delta}_d$ por departamento, en escala logit.}
\label{tab:delta-dep}\\
\hline
 & \multicolumn{2}{c}{E1} & \multicolumn{2}{c}{E2} & \multicolumn{2}{c}{E3} \\
Departamento & M2 & M3 & M2 & M3 & M2 & M3 \\
\hline
\endfirsthead
\hline
Departamento & M2 & M3 & M2 & M3 & M2 & M3 \\
\hline
\endhead
\hline
\endfoot
Antioquia & 0.13 & $-0.28$ & 0.12 & $-0.24$ & $-0.00$ & $-0.24$ \\
Atlántico & 0.77 & 0.79 & 0.70 & 0.77 & 0.46 & 0.73 \\
Bogotá D.C. & 0.60 & 0.94 & 0.34 & 0.99 & $-0.18$ & 0.58 \\
Bolívar & 0.00 & $-0.09$ & 0.13 & $-0.16$ & 0.13 & $-0.05$ \\
Boyacá & $-0.38$ & $-0.22$ & $-0.43$ & $-0.20$ & $-0.36$ & $-0.46$ \\
Caldas & 0.16 & 0.02 & $-0.15$ & $-0.29$ & $-0.29$ & $-0.40$ \\
Caquetá & $-0.34$ & $-0.22$ & $-0.43$ & $-0.10$ & $-0.40$ & $-0.05$ \\
Cauca & $-0.02$ & 0.07 & $-0.22$ & 0.07 & 0.00 & 0.51 \\
Cesar & 0.23 & 0.10 & 0.10 & $-0.18$ & $-0.35$ & $-0.45$ \\
Córdoba & 0.64 & 0.21 & 0.63 & 0.15 & 0.50 & 0.01 \\
Cundinamarca & $-0.54$ & $-0.20$ & $-0.41$ & $-0.07$ & $-0.38$ & $-0.15$ \\
Chocó & $-0.01$ & $-0.03$ & 0.57 & 0.13 & 0.67 & 0.22 \\
Huila & 0.07 & 0.30 & 0.08 & 0.53 & 0.09 & 0.53 \\
La Guajira & 0.15 & 0.20 & 0.69 & 0.36 & 0.41 & 0.27 \\
Magdalena & $-0.24$ & $-0.37$ & $-0.10$ & $-0.43$ & 0.03 & $-0.20$ \\
Meta & $-0.52$ & $-0.17$ & $-0.33$ & 0.04 & $-0.27$ & $-0.16$ \\
Nariño & 0.23 & 0.19 & $-0.07$ & 0.33 & 0.09 & 0.64 \\
N. de Santander & 0.46 & 0.08 & 0.40 & 0.09 & 0.43 & 0.22 \\
Quindío & 0.12 & 0.33 & 0.17 & 0.29 & $-0.23$ & 0.12 \\
Risaralda & 0.04 & 0.06 & 0.14 & 0.08 & 0.08 & 0.21 \\
Santander & $-0.32$ & $-0.47$ & $-0.26$ & $-0.34$ & $-0.24$ & $-0.34$ \\
Sucre & 0.33 & $-0.03$ & 0.38 & $-0.06$ & 0.31 & $-0.05$ \\
Tolima & 0.04 & 0.26 & 0.22 & 0.35 & 0.24 & 0.33 \\
Valle del Cauca & $-0.11$ & $-0.08$ & $-0.53$ & $-0.36$ & $-0.69$ & $-0.35$ \\
Arauca & 0.12 & 0.14 & 0.16 & 0.20 & 0.29 & 0.10 \\
Casanare & 0.04 & 0.34 & 0.08 & 0.41 & 0.07 & 0.05 \\
Putumayo & $-0.15$ & 0.05 & $-0.49$ & $-0.04$ & $-0.41$ & 0.30 \\
San Andrés & $-0.50$ & $-0.18$ & $-0.68$ & $-0.20$ & $-0.23$ & 0.22 \\
Amazonas & $-0.63$ & $-0.39$ & $-0.66$ & $-0.37$ & $-1.17$ & $-0.76$ \\
Guainía & 0.07 & 0.09 & $-0.04$ & 0.02 & 0.48 & 0.44 \\
Guaviare & $-0.07$ & 0.03 & $-0.29$ & $-0.06$ & $-0.07$ & 0.05 \\
Vaupés & 0.05 & $-0.02$ & $-0.27$ & $-0.12$ & $-0.47$ & $-0.48$ \\
Vichada & 0.04 & 0.03 & 0.16 & 0.04 & 0.22 & 0.21 \\
\end{longtable}
\normalsize

%--------------------------------------------------------------------------------------
\section{Coeficientes del Modelo 3 según el rango de la base espacial}
\label{sec:apendice-coef-K}

La Tabla~\ref{tab:coef-K} presenta los coeficientes del Modelo 3, sobre covariables estandarizadas, estimados con los 33 dominios para cada valor de $K$ y $\rho=0.50$.

\begin{table}[!htbp]
\centering
\footnotesize
\caption{Coeficientes estimados del Modelo 3 según el rango $K$ de la base espacial, con $\rho=0.50$.}
\label{tab:coef-K}
\renewcommand{\arraystretch}{1.15}
\begin{tabular}{l l c c c c c}
\hline
Escenario & Coeficiente & $K=3$ & $K=4$ & $K=6$ & $K=8$ & $K=10$ \\
\hline
E1 & Intercepto & $-2.23$ & $-2.17$ & $-2.45$ & $-2.40$ & $-2.57$ \\
E1 & Repitencia escolar & 0.84 & 0.88 & 1.01 & 1.05 & 1.10 \\
E1 & Tasa de matriculación & 0.30 & 0.29 & 0.40 & 0.38 & 0.45 \\
E1 & Infraestructura básica & $-0.27$ & $-0.28$ & $-0.29$ & $-0.29$ & $-0.24$ \\
E1 & Cobertura de energía rural & 0.16 & 0.14 & 0.13 & 0.12 & 0.14 \\
E1 & Mortalidad infantil & 0.18 & 0.09 & 0.07 & 0.07 & 0.08 \\
E1 & Aseguramiento en salud & $-0.25$ & $-0.23$ & $-0.19$ & $-0.14$ & $-0.11$ \\
E1 & Luminosidad nocturna & $-0.21$ & $-0.23$ & $-0.20$ & $-0.22$ & $-0.20$ \\
E2 & Intercepto & $-1.96$ & $-1.98$ & $-2.17$ & $-2.10$ & $-2.14$ \\
E2 & Repitencia escolar & 0.73 & 0.80 & 0.87 & 0.96 & 0.95 \\
E2 & Infraestructura básica & $-0.29$ & $-0.32$ & $-0.35$ & $-0.33$ & $-0.32$ \\
E2 & Aseguramiento en salud & $-0.23$ & $-0.17$ & $-0.13$ & $-0.04$ & $-0.03$ \\
E2 & Luminosidad nocturna & $-0.20$ & $-0.22$ & $-0.18$ & $-0.22$ & $-0.21$ \\
E3 & Intercepto & $-1.69$ & $-1.72$ & $-1.75$ & $-1.78$ & $-1.96$ \\
E3 & Repitencia escolar & 0.72 & 0.76 & 0.80 & 1.07 & 1.23 \\
E3 & Cobertura de energía rural & 0.05 & 0.01 & $-0.01$ & $-0.03$ & $-0.04$ \\
E3 & Aseguramiento en salud & $-0.21$ & $-0.13$ & $-0.07$ & 0.21 & 0.35 \\
E3 & Luminosidad nocturna & $-0.26$ & $-0.27$ & $-0.26$ & $-0.36$ & $-0.31$ \\
\hline
\end{tabular}
\end{table}

% Vaciar las tablas flotantes pendientes (en particular la Tabla A.4)
% antes de comenzar la sección A.7 y su longtable (Tabla A.5).
\clearpage

%--------------------------------------------------------------------------------------

\section{Sensibilidad de los parámetros ante la exclusión de departamentos influyentes}
\label{sec:apendice-influencia-parametros}

La Tabla~\ref{tab:robustez-parametros-completa} presenta los 69 registros de comparación paramétrica correspondientes a las exclusiones individuales de San Andrés y Providencia, Vichada y Vaupés bajo el escenario E2. Se reportan los parámetros originales estimados con 33 departamentos, sus reestimaciones con 32 departamentos, las diferencias calculadas como el valor reestimado menos el original y los cambios relativos porcentuales.

\begingroup
\scriptsize
\setlength{\tabcolsep}{2.8pt}
\renewcommand{\arraystretch}{1.06}

\begin{longtable}{l c l r r r r}

\caption{Sensibilidad completa de los parámetros de M1, M2 y M3 en E2 ante exclusiones individuales.}
\label{tab:robustez-parametros-completa}\\

\hline
Depto. excluido & Modelo & Parámetro & Original & Reest. & Dif. & Cambio relativo (\%) \\
\hline
\endfirsthead

\multicolumn{7}{l}{\textit{Tabla \thetable\ (continuación)}} \\
\hline
Depto. excluido & Modelo & Parámetro & Original & Reest. & Dif. & Cambio relativo (\%) \\
\hline
\endhead

\hline
\endfoot

San Andrés y Prov. & M1 & Intercepto & \num{20.6235} & \num{13.5003} & \num{-7.1232} & \num{-34.5393} \\
 & M1 & Repitencia & \num{380.8131} & \num{428.5059} & \num{47.6928} & \num{12.5240} \\
 & M1 & Infraestructura & \num{-13.1363} & \num{-23.6841} & \num{-10.5478} & \num{-80.2955} \\
 & M1 & Aseguramiento & \num{-32.2045} & \num{-24.2060} & \num{7.9985} & \num{24.8366} \\
 & M1 & Luminosidad & \num{-1.6777} & \num{1.4543} & \num{3.1320} & \num{186.6843} \\
 & M2 & Intercepto & \num{-2.1209} & \num{-2.7798} & \num{-0.6589} & \num{-31.0689} \\
 & M2 & Repitencia & \num{23.3537} & \num{26.2002} & \num{2.8464} & \num{12.1884} \\
 & M2 & Infraestructura & \num{-1.1984} & \num{-1.4126} & \num{-0.2142} & \num{-17.8776} \\
 & M2 & Aseguramiento & \num{-0.9293} & \num{-0.5573} & \num{0.3720} & \num{40.0283} \\
 & M2 & Luminosidad & \num{-0.3175} & \num{-0.0370} & \num{0.2805} & \num{88.3411} \\
 & M3 & Intercepto & \num{-2.7631} & \num{-2.5441} & \num{0.2190} & \num{7.9275} \\
 & M3 & Repitencia & \num{28.3035} & \num{29.0006} & \num{0.6971} & \num{2.4629} \\
 & M3 & Infraestructura & \num{-1.4022} & \num{-1.4560} & \num{-0.0538} & \num{-3.8378} \\
 & M3 & Aseguramiento & \num{-0.9457} & \num{-0.8881} & \num{0.0576} & \num{6.0863} \\
 & M3 & Luminosidad & \num{-0.4015} & \num{-0.3117} & \num{0.0898} & \num{22.3649} \\
 & M3 & $\alpha_{1}$ & \num{0.8800} & \num{0.0000} & \num{-0.8800} & \num{-99.9948} \\
 & M3 & $\alpha_{2}$ & \num{0.6230} & \num{0.6061} & \num{-0.0170} & \num{-2.7242} \\
 & M3 & $\alpha_{3}$ & \num{0.0632} & \num{0.0416} & \num{-0.0216} & \num{-34.1672} \\
 & M3 & $\alpha_{4}$ & \num{-0.3246} & \num{-0.3187} & \num{0.0059} & \num{1.8034} \\
 & M3 & $\alpha_{5}$ & \num{-0.3315} & \num{-0.3445} & \num{-0.0130} & \num{-3.9214} \\
 & M3 & $\alpha_{6}$ & \num{0.0784} & \num{0.0842} & \num{0.0058} & \num{7.4158} \\
 & M2 & $\tau$ & \num{4.5708} & \num{4.4099} & \num{-0.1610} & \num{-3.5215} \\
 & M3 & $\tau$ & \num{2.5584} & \num{2.5869} & \num{0.0285} & \num{1.1151} \\

\hline

Vichada & M1 & Intercepto & \num{20.6235} & \num{2.6850} & \num{-17.9385} & \num{-86.9808} \\
 & M1 & Repitencia & \num{380.8131} & \num{353.6243} & \num{-27.1888} & \num{-7.1397} \\
 & M1 & Infraestructura & \num{-13.1363} & \num{-11.4702} & \num{1.6660} & \num{12.6826} \\
 & M1 & Aseguramiento & \num{-32.2045} & \num{-11.0026} & \num{21.2019} & \num{65.8352} \\
 & M1 & Luminosidad & \num{-1.6777} & \num{-2.2869} & \num{-0.6092} & \num{-36.3128} \\
 & M2 & Intercepto & \num{-2.1209} & \num{-2.9795} & \num{-0.8586} & \num{-40.4830} \\
 & M2 & Repitencia & \num{23.3537} & \num{22.6070} & \num{-0.7467} & \num{-3.1973} \\
 & M2 & Infraestructura & \num{-1.1984} & \num{-1.0860} & \num{0.1124} & \num{9.3799} \\
 & M2 & Aseguramiento & \num{-0.9293} & \num{0.0369} & \num{0.9662} & \num{103.9681} \\
 & M2 & Luminosidad & \num{-0.3175} & \num{-0.3505} & \num{-0.0330} & \num{-10.3835} \\
 & M3 & Intercepto & \num{-2.7631} & \num{-3.0190} & \num{-0.2558} & \num{-9.2589} \\
 & M3 & Repitencia & \num{28.3035} & \num{27.9581} & \num{-0.3454} & \num{-1.2205} \\
 & M3 & Infraestructura & \num{-1.4022} & \num{-1.3636} & \num{0.0386} & \num{2.7506} \\
 & M3 & Aseguramiento & \num{-0.9457} & \num{-0.6193} & \num{0.3264} & \num{34.5160} \\
 & M3 & Luminosidad & \num{-0.4015} & \num{-0.4203} & \num{-0.0188} & \num{-4.6710} \\
 & M3 & $\alpha_{1}$ & \num{0.8800} & \num{0.8326} & \num{-0.0475} & \num{-5.3921} \\
 & M3 & $\alpha_{2}$ & \num{0.6230} & \num{0.6036} & \num{-0.0195} & \num{-3.1263} \\
 & M3 & $\alpha_{3}$ & \num{0.0632} & \num{0.0440} & \num{-0.0192} & \num{-30.3592} \\
 & M3 & $\alpha_{4}$ & \num{-0.3246} & \num{-0.3364} & \num{-0.0118} & \num{-3.6470} \\
 & M3 & $\alpha_{5}$ & \num{-0.3315} & \num{-0.3303} & \num{0.0012} & \num{0.3621} \\
 & M3 & $\alpha_{6}$ & \num{0.0784} & \num{0.0638} & \num{-0.0146} & \num{-18.5861} \\
 & M2 & $\tau$ & \num{4.5708} & \num{4.5556} & \num{-0.0152} & \num{-0.3335} \\
 & M3 & $\tau$ & \num{2.5584} & \num{2.5921} & \num{0.0337} & \num{1.3184} \\

\hline

Vaupés & M1 & Intercepto & \num{20.6235} & \num{14.8384} & \num{-5.7851} & \num{-28.0509} \\
 & M1 & Repitencia & \num{380.8131} & \num{423.7658} & \num{42.9528} & \num{11.2792} \\
 & M1 & Infraestructura & \num{-13.1363} & \num{-9.1819} & \num{3.9543} & \num{30.1024} \\
 & M1 & Aseguramiento & \num{-32.2045} & \num{-33.1503} & \num{-0.9458} & \num{-2.9369} \\
 & M1 & Luminosidad & \num{-1.6777} & \num{-1.5053} & \num{0.1724} & \num{10.2771} \\
 & M2 & Intercepto & \num{-2.1209} & \num{-2.4228} & \num{-0.3020} & \num{-14.2372} \\
 & M2 & Repitencia & \num{23.3537} & \num{25.1107} & \num{1.7570} & \num{7.5233} \\
 & M2 & Infraestructura & \num{-1.1984} & \num{-0.9369} & \num{0.2614} & \num{21.8168} \\
 & M2 & Aseguramiento & \num{-0.9293} & \num{-0.9788} & \num{-0.0494} & \num{-5.3161} \\
 & M2 & Luminosidad & \num{-0.3175} & \num{-0.2879} & \num{0.0296} & \num{9.3323} \\
 & M3 & Intercepto & \num{-2.7631} & \num{-2.8196} & \num{-0.0564} & \num{-2.0423} \\
 & M3 & Repitencia & \num{28.3035} & \num{28.7743} & \num{0.4707} & \num{1.6632} \\
 & M3 & Infraestructura & \num{-1.4022} & \num{-1.2756} & \num{0.1266} & \num{9.0262} \\
 & M3 & Aseguramiento & \num{-0.9457} & \num{-0.9793} & \num{-0.0336} & \num{-3.5535} \\
 & M3 & Luminosidad & \num{-0.4015} & \num{-0.3987} & \num{0.0028} & \num{0.7089} \\
 & M3 & $\alpha_{1}$ & \num{0.8800} & \num{0.7913} & \num{-0.0887} & \num{-10.0754} \\
 & M3 & $\alpha_{2}$ & \num{0.6230} & \num{0.6111} & \num{-0.0119} & \num{-1.9100} \\
 & M3 & $\alpha_{3}$ & \num{0.0632} & \num{0.0651} & \num{0.0018} & \num{2.8600} \\
 & M3 & $\alpha_{4}$ & \num{-0.3246} & \num{-0.3070} & \num{0.0176} & \num{5.4108} \\
 & M3 & $\alpha_{5}$ & \num{-0.3315} & \num{-0.3178} & \num{0.0137} & \num{4.1447} \\
 & M3 & $\alpha_{6}$ & \num{0.0784} & \num{0.0757} & \num{-0.0027} & \num{-3.4337} \\
 & M2 & $\tau$ & \num{4.5708} & \num{4.4799} & \num{-0.0909} & \num{-1.9894} \\
 & M3 & $\tau$ & \num{2.5584} & \num{2.5564} & \num{-0.0019} & \num{-0.0745} \\

\hline
\end{longtable}
\endgroup

\noindent\textit{Nota:} Los coeficientes fijos de M2 y M3 se expresan en una escala común para facilitar la comparación entre ajustes. Los parámetros $\alpha_k$ son coeficientes de la base espacial de rango reducido, no efectos municipales individuales. El cambio relativo procede del archivo de resultados; cuando el valor original es pequeño o cercano a cero, su interpretación porcentual resulta inestable y debe priorizarse la diferencia absoluta. Los cambios de signo se observan en luminosidad nocturna de M1 al excluir San Andrés y Providencia y aseguramiento en salud de M2 al excluir Vichada.